%% latex_185_malik_s.tex
%%
%% The following file is a skeleton file that demonstrates how to implement the IEEEtran.cls class 
%% file in LaTeX. This file is intended for the CMPE185 Technical Writing for Engineers Course in 
%% which the students must write a tutorial aimed towards novice LaTeX users in the LaTex environment.
%%
%% This file is heavily adapted from Michael Shell's bare_adv.tex file made available at 
%% http://www.ieee.org/conferences_events/conferences/publishing/templates.html
%%
%% You will need to rename this file with your information in the following format:
%% latex_185_last name_ first initial.tex
%% ---------------------------------------------------------------------------------------------------

% \documentclass{} precedes the preamble and is typically the first command in any .tex document. Every .tex document should include this command as it defines what kind of document you intend on creating. Modifiers in square brackets [] can be added in between the text ''documentclass'' and the curly brackets {} to modify font size, templates, etc. 

\documentclass[12pt,journal,compsoc]{IEEEtran}

%-----PACKAGES-------------------------------------------------------------------------------------

% Packages include extra commands that allow for additional formatting, ranging from Graphics, Math, to Alignment. The command to include packages will always look similar to \usepackage{} where the package name is within the curly brackets {}. Packages are defined in the preamble, i.e., between the \documentclass{} and \begin{document} commands.

% The following link takes you to a list of additional packages that may not be listed here. http://en.wikibooks.org/wiki/LaTeX/Package_Reference

% TO INCLUDE PACKAGES: 
% - Open the file ''latex_sample_packages.tex'' from the .zip folder.
% - From this file, COPY the code for the package you want to include and PASTE into your own .tex file.
% - Uncomment the package you want to include and load, i.e., remove the ''%'' in front of the \usepackage{package_name}.
 

% Copy package text here:

\usepackage{graphicx}
\usepackage{textcomp}

% This is an example of how a \usepackage{} command should be included. You will need to include more packages to complete this assignment.

%----- The DOCUMENT Environment-------------------------------------------------------------------

% The \begin{document} and \end{document} commands establish the environment for the text of the document. The \begin{} and \end{} commands are used repeatedly in LaTeX to show where an environment begins and ends. The \end{document} command will be the last line of this .tex file.

\begin{document}

% The following commands are self explanatory. Insert your title, author and date into each command's curly bracket. You can include the abstract, paper header and paper footer information in this section then conclude the section with the command \maketitle (shown as the last line at the end of this section).

\title{\LaTeX\ Tutorial}
\author{Shanaya Malik}
% The double backslash \\ is used here to enter a ''Carriage Return'', or a line break. Note that the tilde ~ in between my name used as a ''Nonbreaking Space''. LaTeX will not break a structure at a ~ so this keeps an author's name from being broken across two lines. Also note that I included my name as an example so make sure to only insert your name in this command.

\date{}		% leaving the brackets empty omits the date
% To input the current date, you can type: \date{\today}

% The paper headers
\markboth{\LaTeX\ IEEE Template Tutorial}%
{Moulds \MakeLowercase{\textit{et al.}}: CMPE185}
% The only time the second header will appear is for the odd numbered pages
% after the title page when using the twoside option.
% 
% *** Note that you probably will NOT want to include the author's ***
% *** name in the headers of peer review papers.                   ***
% You can use \ifCLASSOPTIONpeerreview for conditional compilation here if
% you desire.

% The publisher's ID mark at the bottom of the page is less important with
% Computer Society journal papers as those publications place the marks
% outside of the main text columns and, therefore, unlike regular IEEE
% journals, the available text space is not reduced by their presence.
% If you want to put a publisher's ID mark on the page you can do it like
% this:
\IEEEpubid{0000--0000/00\$00.00~\copyright~2007 IEEE}
% or like this to get the Computer Society new two part style.
%\IEEEpubid{\makebox[\columnwidth]{\hfill 0000--0000/00/\$00.00~\copyright~2007 IEEE}%
%\hspace{\columnsep}\makebox[\columnwidth]{Published by the IEEE Computer Society\hfill}}
% Remember, if you use this you must call \IEEEpubidadjcol in the second
% column for its text to clear the IEEEpubid mark (Computer Society jorunal
% papers don't need this extra clearance.)


% use for special paper notices
%\IEEEspecialpapernotice{(Invited Paper)}

% for Computer Society papers, we must declare the abstract and index terms
% PRIOR to the title within the \IEEEcompsoctitleabstractindextext IEEEtran
% command as these need to go into the title area created by \maketitle.
\IEEEcompsoctitleabstractindextext{%
\begin{abstract}
%\boldmath
The abstract goes here.
\end{abstract}
% IEEEtran.cls defaults to using nonbold math in the Abstract.
% This preserves the distinction between vectors and scalars. However,
% if the journal you are submitting to favors bold math in the abstract,
% then you can use LaTeX's standard command \boldmath at the very start
% of the abstract to achieve this. Many IEEE journals frown on math
% in the abstract anyway. In particular, the Computer Society does
% not want either math or citations to appear in the abstract.

% Note that keywords are not normally used for peerreview papers.
\begin{IEEEkeywords}
CMPE185, \LaTeX\ Tutorial, IEEEtran, journal, \LaTeX, paper, template.
\end{IEEEkeywords}}

\maketitle

%----- The SECTION Environment -------------------------------------------------------------------

% To create a section, simply type the command \section{} with the name of your section name inserted into the curly brackets {}. The section's body text follows underneath the \section{} command. 

\section{Introduction}
% Computer Society journal papers do something a tad strange with the very
% first section heading (almost always called "Introduction"). They place it
% ABOVE the main text! IEEEtran.cls currently does not do this for you.
% However, You can achieve this effect by making LaTeX jump through some
% hoops via something like:
%
%\ifCLASSOPTIONcompsoc
%  \noindent\raisebox{2\baselineskip}[0pt][0pt]%
%  {\parbox{\columnwidth}{\section{Introduction}\label{sec:introduction}%
%  \global\everypar=\everypar}}%
%  \vspace{-1\baselineskip}\vspace{-\parskip}\par
%\else
%  \section{Introduction}\label{sec:introduction}\par
%\fi
%
% Admittedly, this is a hack and may well be fragile, but seems to do the
% trick for me. Note the need to keep any \label that may be used right
% after \section in the above as the hack puts \section within a raised box.



% The very first letter is a 2 line initial drop letter followed
% by the rest of the first word in caps (small caps for compsoc).
% 
% form to use if the first word consists of a single letter:
% \IEEEPARstart{A}{demo} file is ....
% 
% form to use if you need the single drop letter followed by
% normal text (unknown if ever used by IEEE):
% \IEEEPARstart{A}{}demo file is ....
% 
% Some journals put the first two words in caps:
% \IEEEPARstart{T}{his demo} file is ....
% 
% Here we have the typical use of a "T" for an initial drop letter
% and "HIS" in caps to complete the first word.

\IEEEPARstart{T}{his} tutorial is designed to introduce novice engineering students to \LaTeX\, a software system for typesetting documents, which is widely used in both academia and the professional engineering world. The tool operates on a principle that allows the writer to focus primarily on the content, using plain text and commands to control formatting.
% You must have at least 2 lines in the paragraph with the drop letter
% (should never be an issue)

% Creating a subsection is similar to creating a section and is used with the command \subsection{}.

\subsection{Why will this tutorial be helpful?}
\LaTeX\ is meant for managing complex mathematical equations, tables, and figures, making it an ideal choice for engineering documents. The tutorial will guide students through the process of creating documents that are not only visually professional but also precisely formatted.

% needed in second column of first page if using \IEEEpubid
%\IEEEpubidadjcol

% Creating a subsubsection:

%\subsubsection{Why should I learn \LaTeX\?}
\subsubsection{Why should I learn \LaTeX{}?}
Familiarity with LaTeX allows students to produce well-structured reports and papers that adhere to the standards required in scientific communication. \LaTeX\ documents maintain consistency across different computing platforms, improving collaboration on projects. 

%---- Additional Features ----
\section{Creating a \LaTeX\ File}
\IEEEPARstart{I}{n} order to create a '.tex' file, it is required to use a text editor, but no specialized software is needed. There are editors, such as Overleaf, which include \LaTeX\ syntax highlighting and integration, and optimize the process.

\subsection{Environments}

Environments in \LaTeX\ determine the formatting of text, and are represented as \texttt{\textbackslash begin} and \texttt{\textbackslash end} commands. It can be used for a variety of purposes, such as itemizing lists; aligning text; or formatting mathematical expressions.

\subsubsection{Begin}

The \texttt{\textbackslash begin\{environment\}} command initiates the environment, where \texttt{environment} is the name of the environment you wish to use.  

\subsubsection{End}

The \texttt{\textbackslash end\{environment\}} command is used to signal the end of an environment. \LaTeX\ will stop applying the special formatting of the environment after this command.  

\subsection{Reserved Characters}

In \LaTeX\, there are characters that have special meanings and are 'reserved' for specific commands and functions. The characters include symbols, such as the backslash `\`, which indicates the start of a command, and the percent `\%`, which comments out text, to name a few. In order to display these characters as part of your text, the character must be escaped using specific commands. 

\subsubsection{Why are these reserved?}

Reserved characters in LaTeX serve specific functions in the typesetting process, such as invoking commands or formatting text. Because of these roles, they cannot be used directly to display their literal symbols without being escaped.

\subsubsection{Explain the functions of these characters: \textbackslash , \textasciitilde , \textbackslash\textbackslash , \%}

\begin{itemize}
    \item \textbackslash : Introduces a command or a LaTeX keyword.
    \item \textasciitilde : Produces a non-breakable space.
    \item \textbackslash\textbackslash : Creates a new line or a line break.
    \item \% : Starts a comment, which makes LaTeX ignore the rest of the text on the line.
\end{itemize}

\subsubsection{What if you want to display these characters?}

To display reserved characters in the text, you need to use a specific set of commands that tells LaTeX to treat them as regular characters and not as commands or special symbols.

\begin{itemize}
    \item To display a backslash, use \texttt{\textbackslash textbackslash}.
    \item To display a tilde, use \texttt{\textbackslash textasciitilde} or place it inside math mode as \texttt{\$ \textasciitilde \$}.
    \item To display double backslashes, use \texttt{\textbackslash textbackslash\textbackslash textbackslash}.
    \item To display the percent sign, use \texttt{\textbackslash textpercent}.
\end{itemize}

%----- Additional Features -----------------------------------------------------------------------

\section{Additional Features}

\subsection{Figures}
% FIGURES:

% Note the FIGURE Environment created by the \begin{figure} and \end{figure} commands.

\begin{figure}[h] 	% There are several different modifiers that can be used in [].
\centering
\includegraphics[width=1.5in]{slug.pdf}
\caption{Sammy the Slug}
\label{fig_slug}
\end{figure}

% You will need to use appropriate file types for figures and will also need to include that image file in the same folder as your .tex file. 

%-------------------------------------------------------------------------------------------------
\subsection{Label, Cite, and Ref Commands}
You will need to be able to define and explain how to use the following commands:
\begin{verbatim} \label{fig_slug} \end{verbatim} as it appears in the figure environment\\
\begin{verbatim} \cite{IEEEhowto:kopka} \end{verbatim} appears like: \cite{IEEEhowto:kopka}\\
\begin{verbatim} \ref{fig_slug} \end{verbatim} appears like: \ref{fig_slug}\\

% IMPORTANT NOTE: In order to assign the correct reference number to each label, you may have to compile your code twice. 

%-------------------------------------------------------------------------------------------------
\subsection{Tables}
An example of a floating table.

% An example of a floating table. Note that, for IEEE style tables, the 
% \caption command should come BEFORE the table. Table text will default to
% \footnotesize as IEEE normally uses this smaller font for tables.
% The \label must come after \caption as always.
%
\begin{table}[h]
%% increase table row spacing, adjust to taste
\renewcommand{\arraystretch}{1.3}
% if using array.sty, it might be a good idea to tweak the value of
%\extrarowheight as needed to properly center the text within the cells
\caption{An Example of a Table}
\label{table_example}
\centering
%% Some packages, such as MDW tools, offer better commands for making tables
%% than the plain LaTeX2e tabular which is used here.
\begin{tabular}{|c||c|}
\hline
One & Two\\
\hline
Three & Four\\
\hline
\end{tabular}
\end{table}


% Note that IEEE does not put floats in the very first column - or typically
% anywhere on the first page for that matter. Also, in-text middle ("here")
% positioning is not used. Most IEEE journals use top floats exclusively.
% However, Computer Society journals sometimes do use bottom floats - bear
% this in mind when choosing appropriate optional arguments for the
% figure/table environments.
% Note that, LaTeX2e, unlike IEEE journals, places footnotes above bottom
% floats. This can be corrected via the \fnbelowfloat command of the
% stfloats package.


\section{Conclusion}
Conclusion goes here.

%----- APPENDICES --------------------------------------------------------------------------------
\appendices
\section{Appendix Title}
Appendix one text goes here.

% you can choose not to have a title for an appendix
% if you want by leaving the argument blank
\section{}
Appendix two text goes here.


%----- ACKNOWLEDGEMENT SECTION -------------------------------------------------------------------
% Explain what the asterisk * does in the next line: 
\section*{Acknowledgements}

The author would like to thank...\\ \\

Reminder: you will need to explain how to include an Acknowledgement Section and then include your own Acknowledgement Section at the end of your own tutorial. Same applies for the References/Bibliography.


%----- BIBLIOGRAPHY ------------------------------------------------------------------------------

% You will need to explain how to include the bibliography section as follows. Explain the environment and how to add new items.
% Including how \ref, \cite and \label should be included here.

% Reminder: you will need to explain how to include the Bibliography Section and then include your own Bibliography at the end of your own tutorial.

\begin{thebibliography}{1}

\bibitem{IEEEhowto:kopka}
H.~Kopka and P.~W. Daly, \emph{A Guide to {\LaTeX}}, 3rd~ed.\hskip 1em plus
  0.5em minus 0.4em\relax Harlow, England: Addison-Wesley, 1999.

\end{thebibliography}

%----- Optional: BIOGRAPHY Section ---------------------------------------------------------------
 
% If you have an EPS/PDF photo (graphicx package needed) extra braces are
% needed around the contents of the optional argument to biography to prevent
% the LaTeX parser from getting confused when it sees the complicated
% \includegraphics command within an optional argument. (You could create
% your own custom macro containing the \includegraphics command to make things
% simpler here.)
%\begin{biography}[{\includegraphics[width=1in,height=1.25in,clip,keepaspectratio]{mshell}}]{Gerald Moulds}
% or if you just want to reserve a space for a photo:

\begin{IEEEbiography}{Gerald Moulds}
Biography text here.
\end{IEEEbiography}

% if you will not have a photo at all:
\begin{IEEEbiographynophoto}{John Doe}
Biography text here.
\end{IEEEbiographynophoto}

% insert where needed to balance the two columns on the last page with
% biographies
%\newpage

\begin{IEEEbiographynophoto}{Jane Doe}
Biography text here.
\end{IEEEbiographynophoto}

% You can push biographies down or up by placing
% a \vfill before or after them. The appropriate
% use of \vfill depends on what kind of text is
% on the last page and whether or not the columns
% are being equalized.

%\vfill

% Can be used to pull up biographies so that the bottom of the last one
% is flush with the other column.
%\enlargethispage{-5in}

\end{document}
